\clearpage
\begingroup
\fancyhead{}
\renewcommand{\headrulewidth}{0pt}
\fancypagestyle{VIIPrinted641}{
  \fancyhead[C]{\parbox[b]{\dimexpr\headwidth-44pt\relax}{\centering\small
    BEMERKUNGEN.\quad THEORIE DER BEWEGUNG DES MONDES.}}
  \fancyhead[R]{\small 641}
}
\SourceLeaf{646}
\needspace{6\baselineskip}
\pdfbookmark[1]{Bemerkungen zur Theorie der Bewegung des Mondes}{moon-bemerkungen}
\begin{center}BEMERKUNGEN ZUR THEORIE DER BEWEGUNG DES MONDES.\end{center}
\small

\par\noindent\hspace*{\parindent}Die vorstehend abgedruckte »Theorie der Bewegung des Mondes«, die auf einem Convolut von
Blättern steht, ist unvollendet geblieben; ausser dem abgedruckten befinden sich im Nachlass noch zwei
frühere weniger ausgeführte Entwürfe. In \textsc{Gauss}’ wissenschaftlichem Tagebuche findet sich unter 1801
August die Notiz: »\textit{Theoriam motus Lunae aggressi sumus}«; die Entstehung des Manuskripts wird man
thatsächlich in diese Zeit setzen müssen; neben der Handschrift und der Darstellungsweise spricht auch der
Umstand für diese frühe Entstehungsperiode, dass die Anomalien vom Apogäum statt vom Perigäum gezählt
sind, vgl. dazu die Ceresstörungen.

\par\noindent\hspace*{\parindent}Das Interesse, das diese \textsc{Gauss}schen Untersuchungen gegenüber den Methoden anderer, auch späterer
Autoren, bieten, concentrirt sich auf die ihm eigene Behandlung der Differentialgleichungen der Mondbewegung.
Die Durchführung der Integration und die Ermittelung der Störungsglieder ist unvollständig
geblieben, nur die Breitenstörungen sind vollständig, und offenbar hat \textsc{Gauss} diese Papiere später nie
wieder zur Hand genommen; auch hat er sich sonst nirgends darüber geäussert, auch nicht in dem 1802
beginnenden Briefwechsel mit \textsc{Olbers}. Im Jahr 1802 erschien auch bereits der dritte Band von \textsc{Laplaces}
Mécanique céleste, enthaltend die Mondtheorie, wodurch \textsc{Gauss}’ begonnene Entwickelungen überholt
wurden. Auch dieser Umstand, sowie der direkte Anschluss an \textsc{Tobias Mayer} spricht für die erwähnte
frühe Abfassungszeit.

\par\noindent\hspace*{\parindent}Die Form, in der sich schliesslich bei \textsc{Gauss} die Störungen darstellen, ist dieselbe wie die \textsc{Plana}sche
(1832), da beide die Divisoren nach Potenzen der Grösse \(n\) (Verhältniss der Umlaufszeiten von Mond und
Sonne) entwickeln.

\par\noindent\hspace*{\parindent}Einige kleine Zusätze und Verbesserungen sind gemacht worden, um die Resultate wenigstens in den
von \textsc{Gauss} zunächst angenommenen Grenzen vollständig und richtig zu geben. So sind bei der Integration
der zweiten Gleichung des dritten Abschnitts, S.~622, die Glieder in \(\varepsilon\sin(2E\pm a)\) hinzugefügt und in
den folgenden Gleichungen beibehalten worden, obwohl sie, wie \textsc{Gauss} bemerkt, dort fortgelassen werden
können, weil sie auch nach den Integrationen nur von der dritten Ordnung sind. Indessen ist ihre Berücksichtigung
für die zweite Annäherung erforderlich, um die Coefficienten der analogen Glieder im Ausdruck
für \(\frac{\mathrm du^2}{\mathrm dv^2}\pi\), S.~634, richtig zu erhalten.

\par\noindent\hspace*{\parindent}Im vierten Abschnitt ist die Differentialgleichung für \(\theta\) bis auf die Glieder 5.~Ordnung einschliesslich
genau; nach der Integration fehlen also naturgemäss die Glieder, welche durch den Divisor \(n\) in die
5.~Ordnung aufrücken, sowie das Glied in \(\sin(2E-\lambda+2a)\), das den Divisor \(nn\) erhält, und dessen
Fehlen \textsc{Gauss} durch einen \(*\) andeutet; es ergibt sich nach \textsc{Plana} zu \(-\frac{9}{32}n\varepsilon\varepsilon\sin(2E-\lambda+2a)\).
Bei der Transformation auf die \textsc{Mayer}sche Form, S.~631---633, sind die Verzeichnisse der hinzutretenden Glieder,
welche übrigens nicht vollständig sind, beim Abdruck fortgelassen worden. Auch ist die Gegenüberstellung
der \textsc{Gauss}schen und \textsc{Mayer}schen Resultate keiner Controllrechnung unterworfen worden.

\clearpage
\SourceLeaf{647}\thispagestyle{VIIPrinted641}
\par\noindent\hspace*{\parindent}Im fünften Abschnitt hat \textsc{Gauss} die Differentialgleichung für \(p\) nur bis zu den Gliedern vierter Ordnung
einschliesslich aufgestellt; bei der Entwickelung des Ausdrucks für \(\sin 2(v-V)\) sind indessen einige
Glieder dritter Ordnung mitgenommen, weil sie bei der Integration kleine Divisoren erhalten; die beiden
Glieder in Klammern sind hier lediglich deswegen hinzugesetzt, weil \textsc{Gauss} bei Aufstellung des darauf
folgenden Ausdrucks sie berücksichtigt hat. Der schliessliche Ausdruck von \(p\) enthält selbstverständlich
die Glieder niederer als 5.~Ordnung nicht, welche bei der Integration aus solchen 5.~Ordnung entstehen;
einzelne dieser Glieder hatte \textsc{Gauss} zwar schon in seinem Manuscript nachträglich hinzugefügt. Da
das aber nur ganz vereinzelt und flüchtig geschehen ist, und auch vielfach noch nicht die richtigen Werthe
gegeben waren, so wurden sie nicht mitabgedruckt; auch sonst fanden sich einige Unrichtigkeiten vor,
die verbessert worden sind; \textsc{Gauss} hat hier offenbar seine Untersuchungen ziemlich plötzlich abgebrochen,
ohne seine Rechnungen bereits geprüft zu haben.
\begin{flushright}\textsc{Brendel}.\end{flushright}
\par\begin{center}\rule{.24\linewidth}{.4pt}\end{center}
\clearpage
\endgroup
